\clearpage

\setcounter{section}{6}

\section{Das freie Elektronengas im Festkörper}
Bei freien Elektronen handelt es sich um Valenzelektronen von Atomen. Wir betrachten mal die delokalisierten Elektronen, von welchen Alkalimetalle wie \ce{Li}, \ce{Na} usw. eins pro Atom und Erdalkalimetalle wie \ce{Be}, \ce{Mg} usw. zwei pro Atom haben. Diese Elektronen bilden ein Gas mit einer hohen Dichte von $\sim \SI{e22}{\centi\meter^{-3}}$ und haben eine mittlere freie Weglänge von $\overline{\lambda} = \SI{1}{\centi \meter}$ bei sehr reinen Metallen und niedrigen Temperaturen.\\
Bei der Untersuchung der Leitfähigkeit verschiedener Stoffe ergibt sich die Frage, wieso genau kondensierte Materie so durchlässig ist für Leitungselektronen. Aufschluss geben die Tatsachen, dass
\begin{enumerate}
	\item die Elektronenwellen von den Ionenrümpfen eines periodischen Gitters nicht abgelenkt werden (Blochfunktion)
	\item ein Elektron nur selten an anderen Elektronen gestreut wird (Pauli-Prinzip) 
\end{enumerate}
Freie Elektronen bilden in kondensierter Materie also das \textbf{Fermi-Gas freier Elektronen}.

\begin{figure}[H]
    \centering
    \incfig{vl25_abbildung_1}
    \label{fig:vl25_abbildung_1}
    \caption{Qualitativer Verlauf des Potentials für \textbf{ein} Kristallelektron im periodischen Gitter der positiven Rümpfe und aller anderen Elektronen.}
\end{figure}

Bei der einfachsten Beschreibung dieses Fermigases handelt es sich um ein Kastenpotential mit unendlich hohen Wänden. Die Elektronen mit Masse $m$ befinden sich in einem Würfel mit der Kantenlänge $L$. Nach der Schrödingergleichung für ein freies Teilchen welches nur kinetische Energie besitzt, also $V(\Vec{r}) = 0$, gilt
\begin{equation}
	\label{eq:II.7.1}
	-\frac{\hbar^2}{2m} \nabla ^2 \psi(\Vec{r}) = E \psi(\Vec{r}).
\end{equation}
Diese wird gelöst mit dem Ansatz
\begin{equation}
	\label{eq:II.7.2}
	\psi(\Vec{r}) = \frac{1}{L^{3 /2}} e^{i\Vec{k}\Vec{r}},
\end{equation}
wobei $\Vec{k} = (k_{x}, k_{y}, k_{z})$. Die Eigenschaften dieser Lösung sind
\begin{enumerate}
	\item Betrag
		$$
		\psi \psi^{*} = \frac{1}{L^3} e^{i\Vec{k}\Vec{r}}e^{-i\Vec{k}\Vec{r}} = \frac{1}{L^3}
		$$ 
	\item Normierung
		$$
		\int_{\text{Würfel}} \psi \psi^{*} \mathrm{d}V = \frac{1}{L^3} \int \mathrm{d}V = 1
		$$ 
	\item Energie 
		\begin{equation}
			\label{eq:II.7.3}
			E = \frac{\hbar ^2 k^2}{2m}
		\end{equation}
\end{enumerate}
Die ebenen Wellen $\psi_{k}(\Vec{r})$ sind Eigenfunktionen des Impulsoperators $\frac{\hbar }{i} \nabla $, d.h.
$$
\frac{\hbar}{i} \nabla \psi_{k} (\Vec{r}) = \hbar \Vec{k} \psi_{k} (\Vec{r})
$$ 
$$
\implies \Vec{p} = \hbar \Vec{k}
$$ 

\paragraph{Periodische Randbedingungen}\,\\ 
$\psi(x+L, y, z) = \psi (x,y,z)\ldots$
$$
k_{x} = \frac{2 \pi n_{x}}{L}, \quad k_{y} = \frac{2\pi n_{y}}{L}, \quad k_{z} = \frac{2\pi n_{z}}{L}
$$ 
mit $n_{x},n_{y}, n_{z}= 0, \pm 1, \pm 2,\ldots$

$$
K_{n}^2 = \left( \frac{2\pi}{L} \right) ^2 \left( n_{x}^2 + n_{y}^2 + n_{z}^2 \right) = \left( \frac{2\pi}{L} \right) ^2 n^2
$$ 
\begin{equation}
	\label{eq:II.7.4}
	\to E_{n}= \frac{\hbar^2}{2m} \left( \frac{2\pi}{L} \right) ^2 n^2 = \frac{\hbar^2 n^2}{2mL^2}
\end{equation}

\paragraph{Fermienergie}\,\\ 
\textcolor{red}{Ein bisschen Text fehlt hier}
\begin{figure}[H]
    \centering
    \def\svgwidth{0.39\columnwidth}
    \import{figures/}{vl25_abbildung_2.pdf_tex}
    \label{fig:vl25_abbildung_2}
\end{figure}

\begin{equation}
	\label{eq:II.7.5}
	n_{x} = 2 \frac{4}{3} \pi n^{*}^{3} = \frac{8}{3} \pi n^{*}^3
\end{equation}

\begin{equation}
	\label{eq:II.7.6}
	NL^3 = \,\text{Zahl der Elektronen im Kristall}
\end{equation}
mit der Zahl der Elektronen pro Kubikzentimeter $N$. Gleichsetzen von \ref{eq:II.7.5} und \ref{eq:II.7.6} liefert
$$
\begin{aligned}
	\frac{8}{3} \pi n_\mathrm{F}^3 &= NL^3 \\
	n_\mathrm{F} &= \left( \frac{3N}{8\pi} \right)^{1 /3} L.
\end{aligned}
$$
Mit der Energie 
$$
E_\mathrm{F} = \frac{\mathrm{h}^2 n_\mathrm{F}^2}{2mL}
$$ 
ergibt sich
\begin{equation}
	\label{eq:II.7.7}
	E_\mathrm{F} = \frac{\hbar ^2}{2m} \left( 3\pi^2 N \right) ^{2 /3}.
\end{equation}

Für $N=\SI{e 22}{\centi\meter^{-3}}$ ergibt sich nach Gleichung \ref{eq:II.7.7} $E_\mathrm{F} = \SI{2}{\electronvolt}$. $E_\mathrm{F}$ ist nur von der Dichte abhängig und entspricht der höchsten Energie für Elektronen bei $T=\SI{0}{\kelvin}$.

\begin{figure}[H]
    \centering
    \def\svgwidth{0.439\columnwidth}
    \import{figures/}{vl25_abbildung_3.pdf_tex}
    \label{fig:vl25_abbildung_3}
\end{figure}
Aus der Fermienergie lassen sich verschiedene andere Größen berechnen, wie 
\begin{itemize}
	\item die Fermi-Geschwindigkeit, mit
		$$
		E_{\mathrm{F}} = \frac{m}{2} v_{\mathrm{F}}^2 \quad \iff \quad v_{\mathrm{F}} = \sqrt{ \frac{2 E_{\mathrm{F}}}{m}} 
		$$ 
	\item die Fermi-Temperatur, mit
		$$
		E_{\mathrm{F}} = k_{\mathrm{B}}T_{\mathrm{F}} \quad \iff \quad T_{\mathrm{F} } = \frac{E_{\mathrm{F}}}{k_{\mathrm{B}}}
		$$ 
\end{itemize}
Zusätzlich berechnet sich der Fermi-Wellenvektor mit
$$
k_{\mathrm{F}} = \left( 3\pi^2 N \right) ^{1 /3}.
$$ 
Es ist wichtig zu beachten, dass es sich bei $T_{\mathrm{F}}$ nicht um die Temperatur des Elektronengases handelt. Diese beträgt nämlich \SI{0}{\kelvin}!
\begin{table}[htpb]
	\centering
	\caption{caption}
	\label{tab:label}
	\begin{tabular}{|c|c|c|c|c|}
		\hline
		Element & $E_{\mathrm{F}}$ in \si{\electronvolt} & $T_{\mathrm{F}}$ in \unit{\kelvin} & $v_{\mathrm{F}}$ in \si{\centi \meter \per \second} & $k_{\mathrm{F}}$ in $\si{\centi \meter ^{-1}}$ \\
		\hline
		\ce{Na} & \num{3,2} & \num{37000} & \num{1,1 e 8} & \num{0,91 e 8} \\
		\ce{Cu} & \num{7,0} & \num{81000} & \num{1,6 e 8} & \num{1,35 e 8} \\
		\hline
	\end{tabular}
\end{table}

\paragraph{Fermi-Fläche}\,\\ 

\begin{figure}[H]
    \centering
    \incfig{vl25_abbildung_4}
    \label{fig:vl25_abbildung_4}
\end{figure}
Elektronen im Gitter: Fermifläche kann komplizierte Gestalt haben falls die Fermikugel über die 1 BZ. hinausragt (später Theorie der Bänder).

\paragraph{Statistik}\,\\ 
$$
f(E) &= \, \text{Verteilungsfunktion} \hat{=} \, \text{Besetzungswahrscheinlichkeit} \\
&= \frac{n(E)\mathrm{d}E}{D(E) \mathrm{d}E} = \frac{\text{Zahl der Teilchen in} E \ldots E + \mathrm{d}E}{\text{Zahl der Zustände in } E \ldots E + \mathrm{d} E}
$$ 

\paragraph{Fermi-Dirac Statistik}\,\\ 
Die Vorraussetzungungen zur Anwendung der Fermi-Dirac Statistik ist
\begin{enumerate}
	\item Teilchen nicht unterscheidbar
	\item Pauli-Prinzip
\end{enumerate}
\begin{equation}
	\label{eq:II.7.8}
	f(E) = \frac{1}{e^{ \frac{E-\mu}{k_{\mathrm{B}}T}} +1}
\end{equation}
Die Fermi-Dirac Verteilung gibt die Wahrscheinlichkeit dafür an, dass ein Zustand der Energie $E$ im thermischen Gleichgewicht besetzt ist. Dabei wird das chemische Potential $\mu=\mu(T)$, wird in einem speziellen Problem so gewählt, dass die Gesmatzahl der Teilchen im System richtig berücksichtigt wird
$$
n = \int D(E) f(E) \mathrm{d} E
$$ 

\begin{figure}[H]
    \centering
    \incfig{vl25_vorlesung_5}
    \label{fig:vl25_vorlesung_5}
\end{figure}

Beachte: $E_{\mathrm{F}} \approx \SI{3}{\electronvolt}$, $k_{\mathrm{B}}T (\SI{300}{\kelvin}) = \frac{1}{40}\si{ \electronvolt}$.

Für den hochenergetischen Ausläufer der Verteilungsfunktion mit $E-\mu \gg k_{\mathrm{B}} T$ 
\begin{equation}
	\label{eq:II.7.9}
	f(E) = e^{ \frac{\mu - E}{k_{\mathrm{B}}T}}
\end{equation}
(Boltzmannverteilung).

\paragraph{Zustandsdichte (freie Elektronengas)}\,\\ 
Die Zustandsdichteberechnet sich mit Gleichung \ref{eq:II.6.2}
\begin{equation}
	\label{eq:II.7.11}
	D(E) \mathrm{d}E = \frac{2 \cdot V}{(2\pi)^3} \int_{E}^{E+\mathrm{d}E} \mathrm{d}^3 k = \frac{2V}{\left( 2\pi \right) ^3} \frac{1}{V} \mathrm{d}E \int_{E=\, \text{konst.}} \frac{\mathrm{d}S_{E}}{v_{\mathrm{g}}}
\end{equation}
der Faktor $2$ kommt durch Spin \textcolor{red}{?} zustande. Mit
$$
v_{g} = \frac{\partial E}{\partial (\hbar k)} = \frac{\hbar k}{m}
$$ 
und
$$
\int_{\text{Kugel}} \mathrm{d}S_{E} = 4 \pi k^2
$$
folgt
$$
D(E) = \frac{2V}{(2\pi)^3 \hbar } \frac{m}{\hbar k} \cdot 4 \pi k^2
$$ 
und somit
\begin{equation}
	\label{eq:II.7.13}
	D(E) \mathrm{d}E = \frac{V}{2\pi^2} \left( \frac{2m}{\hbar ^2} \right)^{3 /2} \sqrt{E} \mathrm{d}E
\end{equation}

\begin{figure}[H]
    \centering
    \incfig{vl25_abbildung_6}
    \label{fig:vl25_abbildung_6}
\end{figure}

Teilchendichte
\begin{equation}
	\label{eq:II.7.14}
	n(E) \mathrm{d}E = D(E) f(E) \mathrm{d}E = \frac{L^3}{2\pi^2} \left( \frac{2m}{\hbar ^2} \right) ^{3 /2} \frac{\sqrt{E} \mathrm{d}E}{e^{ \frac{E-\mu}{k_{\mathrm{B}}T}} + 1}
\end{equation}

\begin{figure}[H]
    \centering
    \incfig{vl25_abbildung_7}
    \label{fig:vl25_abbildung_7}
\end{figure}

\paragraph{Mittlere Energie des Elektronengases}\,\\ 
Die mittlere Energie pro Elektron berechnet sich mit
$$
\overline{E} = \frac{\int_{0}^{\infty} E n(E) \mathrm{d}E}{\int_{0}^{\infty} n(E) \mathrm{d}E} 
$$ 
Für 
\begin{enumerate}
	\item $T=\SI{0}{\kelvin}$ gilt 
		\begin{equation}
			\label{eq:II.7.15}
			\overline{E} = \frac{3}{5} E_{\mathrm{F}}
		\end{equation}
	\item $T>\SI{0}{\kelvin}$ gilt
		\begin{equation}
			\label{eq:II.7.16}
			\overline{E} = \frac{3}{5} E_{\mathrm{F}} \left( 1 + \frac{5\pi^2}{12} \left( \frac{k_{\mathrm{B}}T}{E_{\mathrm{F}}} \right) ^2 \right) 
		\end{equation}
\end{enumerate}
Es gilt bspw. für $T= \SI{300}{\kelvin}$ 
$$
E_{\mathrm{F}} = \SI{5}{\electronvolt}\quad \implies\quad \overline{E} = \frac{3}{5} E_{\mathrm{F}} + \SI{e-4}{\electronvolt}
$$ 
der thermische Energiebeitrag ist also verschwindend gering.
